\section{تحديث نطاق الفصل بعد الدفعة الثانية}

تتجاوز هذه الدفعة الحدود المعلنة في افتتاح الفصل: فقد أُنجز الآن بناء
الموصل والشخصية البدائية، ومجاميع غاوس، وتحويل ثيتا الملتوي، والاستمرار
العالمي والمعادلة الوظيفية للشخصيات البدائية. وتبقى حالة الفصل
\textenglish{DRAFT}؛ إذ لم يثبت بعد عدم انعدام \(L(1,\chi)\)، ولم
تكتب مبرهنة ديريشليه للأوليات في المتتاليات الحسابية.

\section{الشخصيات المستحثة والموصل}

إذا كان \(d\mid q\)، فإن الاختزال بترديد \(d\) يعطي تشاكلًا شاملًا
\[
\pi_{q,d}:(\mathbb Z/q\mathbb Z)^{\times}
\longrightarrow
(\mathbb Z/d\mathbb Z)^{\times}.
\]
فالشمول ينتج من مبرهنة الباقي الصيني: يمكن رفع كل فئة أولية مع \(d\)
إلى عدد يبقى غير منعدم بترديد كل أولي يقسم \(q\).

\begin{definition}
نسمي \(d\mid q\) \emph{ترديدًا مستحثًا} لشخصية
\(\chi\bmod q\) إذا كانت \(\chi\) تساوي \(1\) على نواة
\(\pi_{q,d}\). عندئذ تنزل \(\chi\) نزولًا وحيدًا إلى شخصية على
\((\mathbb Z/d\mathbb Z)^{\times}\).

وتسمى الشخصية \(\chi\bmod q\) \emph{بدائية} إذا لم يكن لها ترديد
مستحث أصغر من \(q\). أما أصغر ترديد مستحث فيسمى \emph{موصلها}.
\end{definition}

إذا نزلت \(\chi\) إلى شخصية \(\chi^{*}\bmod d\)، فإن المساواة بين
الدالتين الممتدتين على الأعداد الصحيحة ليست حرفية عند الأعداد غير الأولية
مع \(q\). والصيغة الصحيحة هي
\[
\chi(n)=\chi^{*}(n)\chi_{0,q}(n),
\]
حيث \(\chi_{0,q}\) الشخصية الرئيسية بترديد \(q\).

\begin{theorem}[الجد البدائي والموصل]
\label{thm:primitive-ancestor-conductor}
\resultid{ANT-THM-07-04}
\provedhere
لكل شخصية ديريشليه \(\chi\bmod q\) يوجد قاسم وحيد \(f\mid q\)،
وشخصية بدائية وحيدة \(\chi^{*}\bmod f\)، بحيث
\[
\chi(n)=\chi^{*}(n)\chi_{0,q}(n)
\qquad(n\in\mathbb Z).
\]
والعدد \(f\) هو موصل \(\chi\).
\end{theorem}

\begin{proof}
نكتب
\[
q=\prod_{p\mid q}p^{\alpha_p}.
\]
وبمبرهنة الباقي الصيني تنحل زمرة الوحدات إلى
\[
(\mathbb Z/q\mathbb Z)^{\times}
\cong
\prod_{p\mid q}(\mathbb Z/p^{\alpha_p}\mathbb Z)^{\times},
\]
وتنحل الشخصية تبعًا لذلك إلى حاصل ضرب شخصيات محلية \(\chi_p\).

لكل \(p\mid q\)، اختر أصغر عدد
\(0\leq\beta_p\leq\alpha_p\) تكون عنده \(\chi_p\) تافهة على نواة
الاختزال
\[
(\mathbb Z/p^{\alpha_p}\mathbb Z)^{\times}
\longrightarrow
(\mathbb Z/p^{\beta_p}\mathbb Z)^{\times}.
\]
ونفسر زمرة الوحدات عند \(p^0=1\) بأنها الزمرة التافهة. وجود
\(\beta_p\) واضح عند \(\beta_p=\alpha_p\)، وصغراه معرفة لأن المجموعة
منتهية. تنزل \(\chi_p\) نزولًا وحيدًا إلى شخصية
\(\chi_p^{*}\bmod p^{\beta_p}\).

ضع
\[
f=\prod_{p\mid q}p^{\beta_p},
\qquad
\chi^{*}=\prod_{p\mid q}\chi_p^{*}.
\]
إذن \(f\mid q\)، وتستحث \(\chi^{*}\) الشخصية \(\chi\). وإذا كانت
\(\chi^{*}\) مستحثة من قاسم صحيح \(d<f\)، فإن أس أحد الأوليات في
\(d\) أصغر من \(\beta_p\)، وهذا يناقض صغر \(\beta_p\). ومن ثم
\(\chi^{*}\) بدائية.

أما التفرد، فلو استحثت شخصية بدائية \(\psi\bmod g\) الشخصية نفسها، فإن
كل أس أولي في \(g\) يعطي مستوى تنزل إليه الشخصية المحلية، فيكون
\(\beta_p\leq v_p(g)\). ولو كانت المتباينة صارمة عند أولي ما لكانت
المركبة المحلية لـ\(\psi\) مستحثة من قوة أصغر، فتفقد \(\psi\)
بدائيتها. إذن \(v_p(g)=\beta_p\) لكل \(p\)، ومن ثم \(g=f\)، ويعطي
شمول خرائط الاختزال تفرد الشخصية النازلة نفسها.
\end{proof}

\begin{proposition}[العوامل المحلية المحذوفة]
\label{prop:induced-l-factorization}
\resultid{ANT-PROP-07-03}
\provedhere
إذا كانت \(\chi\bmod q\) مستحثة من الشخصية البدائية
\(\chi^{*}\bmod f\)، فإن
\[
L(s,\chi)
=
L(s,\chi^{*})
\prod_{\substack{p\mid q\\p\nmid f}}
\left(1-\frac{\chi^{*}(p)}{p^s}\right)
\qquad(\Re(s)>1).
\]
وتستمر الهوية إلى كل مجال تستمر إليه الدالتان.
\end{proposition}

\begin{proof}
نقارن العوامل الأويلرية. إذا كان \(p\nmid q\)، فلدينا
\(\chi(p)=\chi^{*}(p)\). وإذا كان \(p\mid q\)، فالعامل المحلي
لـ\(L(s,\chi)\) يساوي \(1\). أما في \(L(s,\chi^{*})\)، فالعامل
يساوي أصلًا \(1\) عند \(p\mid f\)، ويحتاج إلى الحذف فقط عندما
\(p\mid q\) و\(p\nmid f\). وهذا يعطي الصيغة في مجال التقارب المطلق،
ثم ينقلها مبدأ الهوية إلى مجالات الاستمرار.
\end{proof}

\begin{example}[شخصية مستحثة بترديد \(6\)]
لتكن \(\chi_3\) الشخصية غير الرئيسية بترديد \(3\). الشخصية
\[
\chi(n)=\chi_3(n)\chi_{0,6}(n)
\]
شخصية بترديد \(6\)، وموصلها \(3\). وبما أن \(\chi_3(2)=-1\)، فإن
\[
L(s,\chi)=L(s,\chi_3)(1+2^{-s}).
\]
هذا المثال يبين أن الاستحثاث لا يعني أن الدالتين الممتدتين بالصفر
متساويتان عند جميع الأعداد الصحيحة.
\end{example}

\section{مجاميع غاوس وتحويل فورييه المنتهي}

نكتب
\[
e_q(x)=e^{2\pi i x/q}.
\]
ولشخصية \(\chi\bmod q\) نعرف مجموع غاوس الملتوي
\[
\tau(\chi,n)
=
\sum_{a\bmod q}\chi(a)e_q(an),
\qquad
\tau(\chi)=\tau(\chi,1).
\]

\begin{theorem}[فصل مجموع غاوس للشخصية البدائية]
\label{thm:primitive-gauss-sum}
\resultid{ANT-THM-07-05}
\provedhere
إذا كانت \(\chi\bmod q\) بدائية، فإن لكل \(n\in\mathbb Z\):
\[
\tau(\chi,n)=\overline{\chi(n)}\tau(\chi).
\]
وفوق ذلك
\[
|\tau(\chi)|=\sqrt q,
\qquad
\tau(\chi)\tau(\overline\chi)=\chi(-1)q.
\]
\end{theorem}

\begin{proof}
إذا كان \((n,q)=1\)، فإن التغيير \(b\equiv an\pmod q\) يعطي
\[
\tau(\chi,n)
=
\overline{\chi(n)}
\sum_{b\bmod q}\chi(b)e_q(b)
=
\overline{\chi(n)}\tau(\chi).
\]

ولنفرض أن \((n,q)>1\)، واختر أوليًا \(p\mid(n,q)\). بما أن \(\chi\)
بدائية، فالقاسم \(q/p\) ليس ترديدًا مستحثًا؛ لذلك يوجد عدد \(c\) يحقق
\[
c\equiv1\pmod{q/p},
\qquad
(c,q)=1,
\qquad
\chi(c)\ne1.
\]
ولأن \(p\mid n\)، فإن \(q\mid(c-1)n\)، ومن ثم
\(e_q(acn)=e_q(an)\). وبما أن الضرب بـ\(c\) يبدل الفئات بترديد \(q\)،
نحصل على
\[
\tau(\chi,n)=\chi(c)\tau(\chi,n),
\]
فيكون المجموع صفرًا. وهذا يساوي الطرف الأيمن لأن \(\chi(n)=0\).

نحسب الآن بطريقة بارسيفال المنتهية:
\begin{align*}
\sum_{n\bmod q}|\tau(\chi,n)|^2
&=
\sum_{a,b\bmod q}\chi(a)\overline{\chi(b)}
\sum_{n\bmod q}e_q(n(a-b))\\
&=q\sum_{a\bmod q}|\chi(a)|^2
=q\varphi(q).
\end{align*}
ومن صيغة الفصل يساوي الطرف الأيسر
\(\varphi(q)|\tau(\chi)|^2\)، فينتج
\(|\tau(\chi)|^2=q\).

وأخيرًا
\[
\overline{\tau(\chi)}
=
\sum_{a\bmod q}\overline{\chi(a)}e_q(-a)
=
\chi(-1)\tau(\overline\chi),
\]
لأن \(\chi(-1)=\pm1\). لذلك
\[
\tau(\chi)\tau(\overline\chi)
=
\chi(-1)|\tau(\chi)|^2
=
\chi(-1)q.
\]
\end{proof}

\begin{corollary}[تحويل فورييه المنتهي للشخصية]
\label{cor:finite-fourier-character}
\resultid{ANT-COR-07-03}
\provedhere
إذا كانت \(\chi\bmod q\) بدائية، فإن
\[
\chi(n)
=
\frac{\tau(\chi)}{q}
\sum_{m\bmod q}\overline{\chi(m)}e_q(-mn)
\qquad(n\in\mathbb Z).
\]
\end{corollary}

\begin{proof}
المجموع في الطرف الأيمن هو
\(\tau(\overline\chi,-n)=\chi(-n)\tau(\overline\chi)\). ومن ثم يساوي
الطرف كله
\[
\frac{\tau(\chi)\tau(\overline\chi)}q
\chi(-1)\chi(n)=\chi(n),
\]
باستعمال \(\tau(\chi)\tau(\overline\chi)=\chi(-1)q\).
\end{proof}

\section{صيغة بواسون وتحويل ثيتا الملتوي}

نعتمد اتفاق تحويل فورييه
\[
\widehat f(y)=\int_{-\infty}^{\infty}f(x)e^{-2\pi ixy}\,dx.
\]

\begin{theorem}[صيغة بواسون على فئة حسابية]
\label{thm:poisson-residue-class}
\resultid{ANT-THM-07-06}
\citedresult{\cite{Apostol1976,MontgomeryVaughan2007}}
إذا كانت \(f\) دالة من فضاء شوارتز، و\(q\geq1\)، فإن
\[
\sum_{k\in\mathbb Z}f(r+kq)
=
\frac1q\sum_{m\in\mathbb Z}
\widehat f\!\left(\frac mq\right)e_q(mr).
\]
\end{theorem}

لتكن الآن \(\chi\bmod q\) بدائية، مع \(q>1\)، ونضع
\[
a=a_\chi=
\begin{cases}
0,&\chi(-1)=1,\\
1,&\chi(-1)=-1.
\end{cases}
\]
ونعرف، لكل \(t>0\)،
\[
\theta_\chi(t)
=
\sum_{n\in\mathbb Z}n^a\chi(n)
\exp\!\left(-\frac{\pi n^2t}{q}\right).
\]
لا يوجد حد ثابت عند \(n=0\)، لأن \(q>1\) و\(\chi(0)=0\).

\begin{theorem}[تحويل ثيتا الملتوي]
\label{thm:twisted-theta-transformation}
\resultid{ANT-THM-07-07}
\provedhere
ضع
\[
\varepsilon_\chi
=
\frac{\tau(\chi)}{i^a\sqrt q}.
\]
عندئذ \(|\varepsilon_\chi|=1\)، ولكل \(t>0\):
\[
\theta_\chi(t)
=
\varepsilon_\chi
t^{-a-1/2}
\theta_{\overline\chi}(1/t).
\]
\end{theorem}

\begin{proof}
نطبق صيغة بواسون على
\[
f_{a,t}(x)=x^a e^{-\pi tx^2/q}.
\]
وحساب التحويل الغاوسي يعطي
\[
\widehat f_{a,t}\!\left(\frac mq\right)
=
(-i)^a q^{1/2}t^{-a-1/2}m^a
\exp\!\left(-\frac{\pi m^2}{qt}\right).
\]
نفصل مجموع \(\theta_\chi(t)\) إلى الفئات \(r\bmod q\)، ثم نستعمل
بواسون:
\begin{align*}
\theta_\chi(t)
&=
\sum_{r\bmod q}\chi(r)
\sum_{k\in\mathbb Z}f_{a,t}(r+kq)\\
&=
\frac1q\sum_{m\in\mathbb Z}
\widehat f_{a,t}\!\left(\frac mq\right)
\sum_{r\bmod q}\chi(r)e_q(mr)\\
&=
\frac{(-i)^a\tau(\chi)}{\sqrt q}
 t^{-a-1/2}
\sum_{m\in\mathbb Z}m^a\overline{\chi(m)}
\exp\!\left(-\frac{\pi m^2}{qt}\right).
\end{align*}
وبما أن \((-i)^a=i^{-a}\) عندما \(a\in\{0,1\}\)، نحصل على الصيغة
المعلنة. أما \(|\varepsilon_\chi|=1\) فتنتج من
\(|\tau(\chi)|=\sqrt q\).
\end{proof}

\section{الدالة المكتملة والمعادلة الوظيفية}

للشخصية البدائية \(\chi\bmod q\)، حيث \(q>1\)، نعرف أولًا في
\(\Re(s)>1\):
\[
\Lambda(s,\chi)
=
\left(\frac q\pi\right)^{(s+a)/2}
\Gamma\!\left(\frac{s+a}{2}\right)L(s,\chi).
\]

\begin{theorem}[الاستمرار العالمي والمعادلة الوظيفية]
\label{thm:primitive-l-functional-equation}
\resultid{ANT-THM-07-08}
\provedhere
تمتد \(\Lambda(s,\chi)\) إلى دالة تامة على \(\mathbb C\)، وتحقق
\[
\Lambda(s,\chi)
=
\varepsilon_\chi
\Lambda(1-s,\overline\chi),
\qquad
|\varepsilon_\chi|=1.
\]
وبالتالي تمتد \(L(s,\chi)\) نفسها إلى دالة تامة.
\end{theorem}

\begin{proof}
إذا كانت \(\Re(s)>1\)، فإن جمع الحدين \(n\) و\(-n\) في
\(\theta_\chi\) يعطي
\[
\theta_\chi(t)
=2\sum_{n=1}^{\infty}n^a\chi(n)
 e^{-\pi n^2t/q}.
\]
وبالتكامل حدًا حدًا نحصل على تمثيل ميلين
\[
\Lambda(s,\chi)
=
\frac12\int_0^\infty
\theta_\chi(t)t^{(s+a)/2}\,\frac{dt}{t}.
\]
تتناقص \(\theta_\chi(t)\) أسيًا عندما \(t\to\infty\). ويحول تحويل
ثيتا سلوكها عند \(t\to0^+\) إلى سلوك
\(\theta_{\overline\chi}(1/t)\) عند اللانهاية، ولذلك يتقارب التكامل
محليًا بانتظام لكل \(s\in\mathbb C\)، ويعرف دالة تامة.

وباستخدام تحويل ثيتا ثم تغيير المتغير \(u=1/t\):
\begin{align*}
\Lambda(s,\chi)
&=
\frac{\varepsilon_\chi}{2}
\int_0^\infty
\theta_{\overline\chi}(1/t)
 t^{(s-a-1)/2}\,\frac{dt}{t}\\
&=
\frac{\varepsilon_\chi}{2}
\int_0^\infty
\theta_{\overline\chi}(u)
 u^{(1-s+a)/2}\,\frac{du}{u}\\
&=
\varepsilon_\chi\Lambda(1-s,\overline\chi).
\end{align*}
وأخيرًا
\[
L(s,\chi)
=
\left(\frac q\pi\right)^{-(s+a)/2}
\frac{\Lambda(s,\chi)}{\Gamma((s+a)/2)}.
\]
ولأن \(1/\Gamma\) دالة تامة، يكون الطرف الأيمن تامًا.
\end{proof}

\begin{remark}
الحالة البدائية ذات الموصل \(1\) هي دالة زيتا نفسها، وقد عولج قطبها
ومعادلتها الوظيفية في الفصل السادس. لذلك اشترطنا هنا \(q>1\).
\end{remark}

\begin{corollary}[الاستمرار العالمي للشخصيات غير الرئيسية]
\label{cor:all-nonprincipal-entire}
\resultid{ANT-COR-07-04}
\provedhere
إذا كانت \(\chi\bmod q\) غير رئيسية، فإن \(L(s,\chi)\) دالة تامة.
أما إذا كانت رئيسية فلها قطب بسيط وحيد عند \(s=1\).
\end{corollary}

\begin{proof}
إذا كانت \(\chi\) غير رئيسية، فجدها البدائي \(\chi^{*}\) غير رئيسي
وموصله أكبر من \(1\). الدالة \(L(s,\chi^{*})\) تامة بالمبرهنة السابقة،
وتضربها صيغة العوامل المحلية المحذوفة في جداء منته من دوال تامة. أما
الحالة الرئيسية فقد أثبتت في القضية
\ref{prop:principal-character-zeta-factorization}.
\end{proof}

\section{الأصفار البديهية والتناظرات}

\begin{corollary}[الأصفار البديهية والشريط الحرج]
\label{cor:dirichlet-l-trivial-zeros}
\resultid{ANT-COR-07-05}
\provedhere
لتكن \(\chi\bmod q\) بدائية و\(q>1\)، وليكن
\(a\in\{0,1\}\) محددًا بالزوجية. عندئذ
\[
L(-a-2n,\chi)=0
\qquad(n=0,1,2,\ldots).
\]
وهذه الأصفار بسيطة في الحالتين الآتيتين:
\begin{itemize}
  \item \(a=1\) ولكل \(n\geq0\).
  \item \(a=0\) ولكل \(n\geq1\).
\end{itemize}
أما الصفر عند \(s=0\) للشخصية الزوجية فهو بسيط إذا وفقط إذا
\(L(1,\overline\chi)\ne0\)، وهي حقيقة ستثبت في دفعة عدم الانعدام.

وكل صفر آخر لـ\(L(s,\chi)\) يقع في الشريط الحرج المغلق
\[
0\leq\Re(s)\leq1.
\]
\end{corollary}

\begin{proof}
لدالة غاما أقطاب بسيطة عندما
\((s+a)/2=0,-1,-2,\ldots\). وبما أن \(\Lambda(s,\chi)\) تامة، يجب أن
تلغي \(L(s,\chi)\) هذه الأقطاب، فتظهر الأصفار المعلنة.

إذا كان \(1-(-a-2n)>1\)، فإن المعادلة الوظيفية تربط قيمة
\(\Lambda(-a-2n,\chi)\) بقيمة عند نصف المستوى \(\Re(s)>1\)، حيث
لا تنعدم \(L\). لذلك تكون قيمة الدالة المكتملة هناك غير صفرية، ويلزم أن
يلغي \(L\) قطب غاما البسيط بصفر بسيط بالضبط. والاستثناء الوحيد هو
\(a=0,n=0\)، إذ يقابل \(s=1\)، فتكون البساطة مكافئة لعدم انعدام
\(L(1,\overline\chi)\).

ولا توجد أصفار عندما \(\Re(s)>1\) بالمنتج الأويلري. وإذا كان
\(\Re(s)<0\) ولم يكن \(s\) أحد المواضع البديهية، فإن عامل غاما منته
وغير صفري؛ وأي صفر هناك ينتقل بالمعادلة الوظيفية إلى صفر في
\(\Re(1-s)>1\)، وهذا محال.
\end{proof}

\begin{proposition}[تناظرات الأصفار]
\label{prop:dirichlet-l-zero-symmetry}
\resultid{ANT-PROP-07-04}
\provedhere
لدينا
\[
\overline{\Lambda(\overline s,\chi)}
=
\Lambda(s,\overline\chi).
\]
فإذا كان \(\rho\) صفرًا للدالة المكتملة الخاصة بـ\(\chi\)، فإن
\[
1-\rho
\quad\text{صفر لـ}\quad\Lambda(s,\overline\chi),
\]
وكذلك \(1-\overline\rho\) صفر للدالة المكتملة الخاصة بـ\(\chi\).
وإذا كانت \(\chi\) حقيقية، ظهرت الأصفار غير الحقيقية في الرباعيات
\[
\rho,\ \overline\rho,\ 1-\rho,\ 1-\overline\rho.
\]
\end{proposition}

\begin{proof}
تثبت علاقة المرافق أولًا من سلسلة ديريشليه في \(\Re(s)>1\)، ثم تمتد
بمبدأ الهوية. وتنتج بقية العلاقات مباشرة من المعادلة الوظيفية ومن أخذ
المرافق المركب.
\end{proof}

\begin{remark}[فرضية ريمان المعممة]
تقول فرضية ريمان المعممة لدوال ديريشليه إن الأصفار غير البديهية تحقق
\(\Re(\rho)=1/2\). لم يثبت هذا الادعاء، ولا تقدم البراهين السابقة دليلًا
عليه؛ فهي تثبت فقط التناظر والحصر في الشريط الحرج.
\end{remark}

\section{مثال: دالة بيتا لديريشليه}

للشخصية \(\chi_4\) البدائية الفردية بترديد \(4\) نجد
\[
\tau(\chi_4)=2i,
\qquad
\varepsilon_{\chi_4}
=
\frac{2i}{i\sqrt4}=1.
\]
ومن ثم تحقق الدالة المكتملة
\[
\Lambda(s,\chi_4)
=
\left(\frac4\pi\right)^{(s+1)/2}
\Gamma\!\left(\frac{s+1}{2}\right)L(s,\chi_4)
\]
المعادلة
\[
\Lambda(s,\chi_4)=\Lambda(1-s,\chi_4).
\]
وتقع أصفارها البديهية عند
\(-1,-3,-5,\ldots\)، وكلها بسيطة.

\section{تمارين الدفعة الثانية}

\begin{enumerate}
  \item أثبت مباشرة شمول خريطة الاختزال
  \((\mathbb Z/q\mathbb Z)^{\times}\to
  (\mathbb Z/d\mathbb Z)^{\times}\) عندما \(d\mid q\).
  \item احسب موصل كل شخصية بترديد \(8\)، وحدد البدائية منها.
  \item تحقق من صيغة العوامل المحلية لشخصية مستحثة من الترديد \(4\) إلى
  الترديد \(12\).
  \item أثبت أن فصل مجموع غاوس لكل \(n\) يستلزم بدائية الشخصية.
  \item اشتق تحويل فورييه المنتهي من تعامد الدوال الأسية بدل استعمال
  حاصل ضرب مجموعي غاوس.
  \item تحقق يدويًا من \(\tau(\chi_4)=2i\).
  \item أعد حساب التحويل الفورييري للدالة
  \(x e^{-\pi tx^2/q}\) بتفاضل تحويل الغاوسي.
  \item استخرج من المعادلة المكتملة الصيغة غير المكتملة الواردة عادةً
  بدلالة \(\Gamma(s)\) ومجموع غاوس.
  \item بيّن أن رتبة صفر \(L(s,\chi)\) عند \(s=0\) لشخصية بدائية زوجية
  تساوي \(1+\operatorname{ord}_{s=1}L(s,\overline\chi)\).
\end{enumerate}

\section{ما بقي قبل اعتماد الفصل}

تبقى البوابة البرهانية التالية منفصلة:

\begin{itemize}
  \item إثبات عدم انعدام \(L(1,\chi)\) لكل شخصية غير رئيسية.
  \item فصل حجة الشخصيات المركبة عن الحالة الحقيقية حيث يلزم تجنب أي
  استدلال دائري.
  \item اشتقاق مبرهنة ديريشليه للأوليات في المتتاليات الحسابية بعد إغلاق
  بوابة عدم الانعدام.
  \item مراجعة ثانية مستقلة للنص كاملًا قبل رفعه إلى
  \textenglish{REVIEWED}.
\end{itemize}
